\chapter{Files and storage}
\label{ch:files}

\section{Where things go}\label{where-things-go}

Clusters have several storage areas with different rules. Learn yours before you create anything big.

\textbf{Kelvin2 only}, according to the \href{https://ni-hpc.github.io/nihpc-documentation/Quickstart\%20Guide/}{quick-start guide}\footnote{\texttt{https://ni-hpc.github.io/nihpc-documentation/Quickstart\%20Guide/}}:

{\def\LTcaptype{none} % do not increment counter
\begin{longtable}[]{@{}
  >{\raggedright\arraybackslash}p{(\linewidth - 6\tabcolsep) * \real{0.2500}}
  >{\raggedright\arraybackslash}p{(\linewidth - 6\tabcolsep) * \real{0.2500}}
  >{\raggedright\arraybackslash}p{(\linewidth - 6\tabcolsep) * \real{0.2500}}
  >{\raggedright\arraybackslash}p{(\linewidth - 6\tabcolsep) * \real{0.2500}}@{}}
\toprule\noalign{}
\begin{minipage}[b]{\linewidth}\raggedright
Area
\end{minipage} & \begin{minipage}[b]{\linewidth}\raggedright
Path
\end{minipage} & \begin{minipage}[b]{\linewidth}\raggedright
Limit
\end{minipage} & \begin{minipage}[b]{\linewidth}\raggedright
Deletion
\end{minipage} \\
\midrule\noalign{}
\endhead
\bottomrule\noalign{}
\endlastfoot
Home & \texttt{/users/\$USER} & 50 GB and 100k files & None automatic \\
Scratch & \texttt{/mnt/scratch2/users/\$USER} & No per-user quota & \textbf{After 90 days without access} \\
Temp & \texttt{/tmp} on each node & The node\textquotesingle s disk & When the session ends \\
McClayRDS & Project paths & By project & QUB projects only \\
\end{longtable}
}

\texttt{quota\ -s} shows your actual allocation.

The key rule: \textbf{scratch is not a backup and not an archive.} Keep a copy of results that matter on your institution\textquotesingle s storage, and ask your support team what, if anything, is backed up.

\section{A tidy project}\label{a-tidy-project}

\begin{lstlisting}[language=bash]
# cluster

# Braces expand: this makes all four folders
mkdir -p ~/project/{code,jobs,logs,envs}

# Big data lives on scratch
mkdir -p /mnt/scratch2/users/$USER/project/{data,outputs,tmp}

# A shortcut, so ~/project/scratch stands for the long path
ln -s /mnt/scratch2/users/$USER/project ~/project/scratch
\end{lstlisting}

\section{Moving files with scp}\label{moving-files-with-scp}

\texttt{scp\ source\ destination} copies securely. A remote location is written \texttt{kelvin:path}, using your shortcut. \textbf{Run it from your laptop}, in a separate terminal from your cluster session.

\begin{lstlisting}[language=powershell]
# laptop

# Upload a file. Quotes are needed because of the space in the path
scp "D:\My Work\data.csv" kelvin:~/project/scratch/data/

# Upload a folder: -r means recursive
scp -r "D:\My Work\code" kelvin:~/project/

# Download: the source is now the remote one
scp kelvin:~/project/scratch/outputs/results.tgz "D:\My Work"
\end{lstlisting}

Without the shortcut, the port needs a \textbf{capital} \texttt{-P}, unlike \texttt{ssh}, which uses a lower-case \texttt{-p}:

\begin{lstlisting}[language=powershell]
# laptop
scp -P 55890 -i $HOME\.ssh\my-kelvin-key data.csv <username>@login.kelvin.alces.network:~/project/
\end{lstlisting}

\textbf{A Windows trap.} Do not end a quoted folder with a backslash. In \texttt{"D:\textbackslash{}My\ Work\textbackslash{}"}, the \texttt{\textbackslash{}"} is read as a literal quote mark, and you get a confusing "No such file or directory". Write \texttt{"D:\textbackslash{}My\ Work"}.

\section{Many small files: pack first}\label{many-small-files-pack-first}

Sending thousands of small files one by one is slow, and hard on shared storage. Pack, send, unpack:

\begin{lstlisting}[language=powershell]
# laptop
# c create, z compress, f file name. --exclude leaves out unwanted folders
tar --exclude=.git -czf code.tgz code
scp code.tgz kelvin:~/project/
\end{lstlisting}

\begin{lstlisting}[language=bash]
# cluster
# x extract. && means each step runs only if the last succeeded
cd ~/project && tar xzf code.tgz && rm code.tgz
\end{lstlisting}

Coming back, pack only what you need, such as summaries rather than every raw output.

\section{Large transfers}\label{large-transfers}

\textbf{Kelvin2 only.} For large amounts of data, use a data movement node rather than a login node: \texttt{dm1.kelvin.alces.network} or \texttt{dm2.kelvin.alces.network}. Tools such as \texttt{rsync} can resume an interrupted transfer. MobaXterm and WinSCP offer drag and drop.

\section{Files that arrive read-only}\label{files-that-arrive-read-only}

Copied files can keep restrictive permissions, leading to \texttt{Permission\ denied} when you edit them. \texttt{chmod\ -R\ u+w\ \textasciitilde{}/project/code} gives you write permission on everything in that folder.

